\section{结论}
本文交付证据约束的技能判别、最小修复与条件路径实现，并以构造诊断、真实社区CLI和外部存档轨迹检验其适用范围。
判别器在机械构造留出检出140/140个已知故障，不等于自然技能总体质量准确率。
本轮355次隔离CLI执行证实：单样本画像异常与阈值被改写两处问题可由单行补丁修复，目标检查分别41/49、42/49提高至49/49；
默认阈值检查未退化，空值语义边界仍保留。

网络侧，前一阶段18节点ECP在三档预算平均70.83\%，相对人工及全局风险仅增加4.17个百分点，宽预算持平且取证总成本更高。
新增自然工具联合回放由固定原流程38/48提高至修复后48/48，但修复后的简单路径同样通过，不能把收益全归给复杂路由。
外部DeepSeek评审因余额不足中断，86个完整配对案例未建立准确性改进，且未覆盖有产物可查的情况。
因此当前可信成果是可核对的反证、范围受限的修复与执行记录，以及未成立优势的实际证据；
通用语义质量模型、完整外部产物验证和大规模自然网络优势仍未建立。
